\begin{center}
\LARGE
A Polynomial System from Differential Equations 
\end{center}

\begin{center}
\Large
Dongming Wang
\end{center}

\noindent
Date:  July 19th (Friday)\\
Time:  15:00-15:30\\

\begin{center}
\large
Abstract
\end{center}

For a class of cubic differential systems the necessary and sufficient conditions for the origin to be
a center were derived by I. S. Kukles in 1944. These conditions were discovered to be incomplete
in the later 1980s with the aid of elimination methods and computer algebra tools. Since then there
have been several attempts to establish the complete conditions without success. The problem of
completing the conditions can be reduced partially to decomposing a large system of polynomial
equations, which can be done in principle by existing algorithms of polynomial elimination.
However, the occurring polynomials are too large in terms of degree and number of terms to be
manageable. Various techniques and software tools have been tried both automatically and
interactively, but the complete conditions have never been obtained.

In this talk, we explain the initial problem, present the polynomial system generated by a program
from Kukles' differential systems, report some of our experiments, and point out the major
difficulties encountered in dealing with the system. The generated polynomial system and some of
its subsystems will be made available electronically. We propose the problem of solving these
systems as an open challenge for testing elimination algorithms and their implementations and call
publicly for solutions from algorithm and software developers. 

